% Source: Lecture8.pdf, pages 1-21 (typed handout). Title on the source: "Lecture 8: Further
% Properties of Potts Model Partition Functions and Connections with Graph-Theoretic Polynomials".
\lecture{8}
\subsection{Some preliminaries from graph theory}

A graph $G$ is defined by its set of vertices (sites, in physics terminology), $V = \{v_1,\ldots,v_{n(G)}\}$ and its set of
edges (bonds in physics terminology), $E$, so we write $G = (V,E)$. The total numbers of vertices and edges in $G$ are
$n(G)$ and $e(G)$, respectively.

\paragraph{Def.} A graph $G$ is \emph{connected} if, starting at a given vertex $v\in V$, there exists a continuous path along
the edges that enables one to move to any other vertex in $V$. For a general graph, we denote $k(G)$ as the number of
connected components of $G$. An example of a connected graph is the line graph $L_3$ consisting of three vertices $v_i$,
$i = 1,2,3$ with an edge $e_{12}$ connecting $v_1$ and $v_2$, and an edge $e_{23}$ connecting $v_2$ and $v_3$. An example
of a disconnected graph is obtained by deleting either of these two edges; the resultant graph has $k = 2$ components. If
one deletes both edges, the resultant graph has $k = 3$ components.

\paragraph{Def.} The \emph{empty} (null) graph $E_n$ is the graph consisting of $n$ vertices with no edges.

\paragraph{Def.} The \emph{complete graph} $K_n$ is the graph consisting of $n$ vertices such that there is a (single) edge
$e_{ij}$ connecting each distinct pair of vertices $v_i$ and $v_j$.

\paragraph{Def.} A \emph{bridge} on a graph is an edge such that deleting it leads to an increase in the number of components
in the graph. For example, consider the line graph $L_3$ with three vertices. This is a connected graph, so $k(L_3) = 1$.
Both of its two edges are bridges, since deleting either one produces a graph with $k = 2$ connected components.

\paragraph{Def.} A \emph{loop} is an edge that extends out from a given vertex and loops back to connect to the same vertex.
In physics we will usually restrict to loopless graphs.

\paragraph{Def.} Two vertices in $G$ are said to be \emph{adjacent} if they are connected by at least one edge. In general,
two adjacent vertices may be connected by more than one edge, although on a regular lattice one generally restricts to the
minimal case in which adjacent vertices are connected by only one edge.

\paragraph{Def.} The \emph{degree} $\Delta_v$ of a vertex $v$ in $G$ is equal to the number of edges that connect to it.

\paragraph{Def.} A graph $G$ is said to be \emph{$\Delta$-regular} if all of its vertices have the same degree, $\Delta$. A 1D
lattice with periodic boundary conditions (BC) is a 2-regular graph. A 2D square lattice with toroidal boundary conditions
is a 4-regular graph, etc. On a $\Delta$-regular lattice graph $\Lambda$, $\Delta_\Lambda$ is the coordination number of the
lattice. For a $\Delta$-regular lattice graph,
\begin{equation}
  e(G) = \frac\Delta2n(G).
\end{equation}
One can also define an effective degree $\Delta_{\text{eff}}$ of a graph as
\begin{equation}
  \Delta_{\text{eff}} = \lim_{n(G)\to\infty}\frac{2e(G)}{n(G)}.
\end{equation}
For example, consider a 1D lattice with free boundary conditions; then the interior vertices have degree $\Delta = 2$, but
the end vertices have degree $\Delta = 1$, so this is not a $\Delta$-regular graph. However, in the limit as
$n(G)\to\infty$, $\Delta_{\text{eff}} = 2$. Similarly, consider a rectangular $L_x\times L_y$ section of a 2D square lattice
with free boundary conditions. Then the interior vertices have degree 4, the other vertices on the four sides have degree 3,
and the corner vertices have degree 2. Again, in the thermodynamic limit $L_x\to\infty$ and $L_y\to\infty$ with finite
nonzero aspect ratio $L_y/L_x$, $\Delta_{\text{eff}} = 4$. This is thus the same as if one used boundary conditions that are
periodic in one direction (sometimes called cylindrical) or in both directions (toroidal BC). Similar comments apply for
higher-dimensional lattices.

\paragraph{Def.} The graph $C_n$ with $n$ vertices is called the \emph{circuit graph}. It is equivalent to a 1D lattice with
$n$ vertices and PBC.

\paragraph{Def.} A \emph{cycle} (circuit) in $G$ is a path on adjacent edges that forms a circuit. The number of linearly
independent cycles on $G$ is denoted $c(G)$. For example, consider the circuit graph $C_4$ with vertex set
$\{v_1,v_2,v_3,v_4\}$. We can think of this as a square with the lower right-hand vertex denoted $v_1$, the upper right-hand
vertex denoted $v_2$, etc. Now add one diagonal edge (bond) joining the upper right-hand vertex $v_2$ to the lower left-hand
vertex $v_4$. The resultant graph $G$ has $n(G) = 4$, edge set $E = \{e_{12},e_{23},e_{34},e_{41},e_{24}\}$ and
$e(G) = 5$. There are three cyclic paths that one can draw: (i) around the outside of the square; (ii) around the upper
left-hand triangle; and (iii) around the lower right-hand triangle. However, only two of these are linearly independent, so
$c(G) = 2$. (Recall similar considerations when using Kirchhoff's laws in electrical circuits; if this were a circuit with
resistors, one would only write down two linearly independent Kirchhoff equations.) A relevant relation is
\begin{equation}
  c(G) = e(G) + k(G) - n(G);
\end{equation}
for example, for the graph above, $n(G) = 4$, $e(G) = 5$, and $k(G) = 1$, so $c(G) = 5+1-4 = 2$.

\paragraph{Def.} A \emph{tree graph} is a connected graph with no cycles. Although the tree graphs with $n = 1,2,3$ are unique,
for $n\ge 4$, there is more than one tree graph. For example, for $n = 4$, there are two tree graphs, namely the line graph
with 4 vertices and the star graph with one central vertex and three outer vertices. We sometimes use the generic symbol $T_n$
for an $n$-vertex tree graph.

\paragraph{Def.} Given a graph $G = (V,E)$, a \emph{spanning subgraph} $G' = (V',E')$ is a graph containing all of the vertices
of $G$, so $V' = V$, and a subset of the edges of $G$, i.e.\ $E'\subseteq E$.

\paragraph{Def.} A graph is \emph{planar} if it can be embedded in a plane with no intersecting edges.

\paragraph{Def.} Given a planar graph $G$, the \emph{planar dual} of $G$, denoted $G^*$, is obtained by mapping each face in $G$
to a vertex in $G^*$, each vertex in $G$ to a face in $G^*$, and then connecting the vertices of $G^*$ with edges. A planar
graph $G$ is self-dual if $G^* = G$.

\subsection{Cluster formula for the Potts model partition function}

Let us consider the Potts model on a general graph $G = (V,E)$. Formulated in the context of statistical mechanics, this
model involves dynamical variables (classical spins) $\sigma_i$ at each vertex $v_i\in V$ that can take on any of the values
in the set $\{1,\ldots,q\}$, with an interaction energy between the spins on adjacent vertices
$\Ham_{\text{int}} = -J\sum_{e_{ij}}\delta_{\sigma_i\sigma_j}$. For simplicity, let us take the external magnetic field to be
zero. Then $\Ham = \Ham_{\text{int}}$ and $-\beta\Ham = K\sum_{e_{ij}}\delta_{\sigma_i\sigma_j}$, where $K = \beta J$ and
$\beta = 1/(\kB T)$, with $T$ the temperature. In thermal equilibrium at a temperature $T$, the partition function is
\begin{equation}
  Z = \sum_{\{\sigma_i\}}e^{-\beta\Ham} = \sum_{\{\sigma_i\}}\exp\left[K\sum_{e_{ij}}\delta_{\sigma_i\sigma_j}\right].
\end{equation}
Let $v_P = e^K-1$ and, for notational simplicity, suppress the subscript $P$, writing this as $v$. Recall the identity
\begin{equation}
  e^{K\delta_{\sigma_i\sigma_j}} = 1 + \left(e^K-1\right)\delta_{\sigma_i\sigma_j} = 1 + v\delta_{\sigma_i\sigma_j}.
\end{equation}
Using this identity, we can write $Z$ as
\begin{equation}
  Z = \sum_{\{\sigma_i\}}\left[\prod_{e_{ij}}\left(1+v\delta_{\sigma_i\sigma_j}\right)\right].
\end{equation}
As above, let $G'$ denote a spanning subgraph of $G$. Then one can rewrite $Z$ as
\begin{equation}\label{eq:cluster}
  Z(G,q,v) = \sum_{G'\subseteq G}q^{k(G')}v^{e(G')},
\end{equation}
where $k(G')$ and $e(G')$ are the number of connected components and edges in $G'$. This is often called the
Fortuin--Kasteleyn cluster formula for the Potts model partition function by physicists, although an equivalent relation had
been derived earlier by mathematicians. This cluster formula is very useful, since it expresses $Z$ in a graph-theoretic
manner, as a sum over contributions from spanning subgraphs without explicit reference to the sum over spin configurations.

If $G$ is the disjoint union $G = G_1\cup G_2$, then $Z(G,q,v) = Z(G_1,q,v)Z(G_2,q,v)$ so, without loss of generality, we
take $G$ to be connected. (But some $G'$ subgraphs are disconnected.)

Since $k(G')\ge 1$ and $e(G')\ge 0$, the cluster representation of $Z$ shows that $Z(G,q,v)$ is a polynomial in $q$ and $v$.

\paragraph{Corollary.} Since $k(G')\ge 1$, the minimum degree of $Z$ as a polynomial in $q$ is 1, i.e.\ $Z(G,q,v)$ always
contains an overall factor of $q$. Since one of the spanning subgraphs is the empty graph with no edges, the minimum degree of
$Z$ as a polynomial in $v$ is 0.

The cluster formula allows one to formally generalize $q$ from nonnegative integers $q = 1,2,3,\ldots$ to nonnegative real
numbers. For the case of ferromagnetic spin--spin couplings, $v\ge 0$, the resultant $Z$ is nonnegative, and is positive if
$q\ne 0$. Physically, since $\Ham$ is real, the Boltzmann weights $e^{-\beta\Ham}$ are nonnegative. One refers to this situation
by saying that a model has a ``Gibbs measure''. So (aside from the marginal case $q = 0$), the generalization above in the case
of the FM sign of the coupling, $J>0$, yields a Gibbs measure.

If one is interested in zeros of $Z$ in the variable $q$, then one naturally generalizes $q$ to complex numbers, but for most of
our discussion, the generalization to real nonnegative $q$ will be sufficient.

\subsubsection*{Proof of the cluster formula}

The proof of the cluster formula is based on the 1--1 correspondence between the terms in
$Z = \sum_{\{\sigma_i\}}\prod_{e_{ij}}\left(1+v\delta_{\sigma_i\sigma_j}\right)$ and spanning subgraphs $G'\subseteq G$. A sketch of
the proof, using, as illustration, $G = C_3$, where $C_n$ is the circuit graph with $n$ vertices. Here
\begin{equation}
  Z(C_3,q,v) = \sum_{\{\sigma_i\}}\prod_{e_{ij}}\left(1+v\delta_{\sigma_i\sigma_j}\right)
  = \sum_{\{\sigma_i\}}\left(1+v\delta_{\sigma_1\sigma_2}\right)\left(1+v\delta_{\sigma_2\sigma_3}\right)\left(1+v\delta_{\sigma_3\sigma_1}\right).
\end{equation}
There are four types of terms that contribute (in correspondence with $\sum_{G'\subseteq G}q^{k(G')}v^{e(G')}$):
\begin{itemize}
  \item the term 1, corresponding to the $G'$ with three disjoint vertices, no edges (so $e(G') = 0$, $k(G') = 3$), for which
    the sum over the $\sigma_i$'s yields $q^3$, i.e.\ one can choose the $\sigma$'s independently in $q^3$ ways;
  \item $v\left(\delta_{\sigma_1\sigma_2}+\delta_{\sigma_2\sigma_3}+\delta_{\sigma_3\sigma_1}\right)$, corresponding to the three
    $G'$s with one edge and one disjoint vertex, so $e(G') = 1$, $k(G') = 2$. The first of these terms contributes if
    $\sigma_1 = \sigma_2$; here one chooses $\sigma_1 = \sigma_2$ in any of $q$ ways, and then $\sigma_3$ independently in any
    of $q$ ways, for a total of $q^2$ ways; similarly for the other two terms, so these terms contribute $3q^2v$;
  \item $v^2\left(\delta_{\sigma_1\sigma_2}\delta_{\sigma_2\sigma_3}+\delta_{\sigma_2\sigma_3}\delta_{\sigma_3\sigma_1}
    +\delta_{\sigma_3\sigma_1}\delta_{\sigma_1\sigma_2}\right)$, corresponding to $G'$s with two edges, so $e(G') = 2$,
    $k(G') = 1$. For the first term to contribute, all of the $\sigma$'s must be equal, and there are $q$ ways of choosing
    them; similarly for the other two terms, so the total contribution is then $3qv^2$;
  \item $v^3\left(\delta_{\sigma_1\sigma_2}\delta_{\sigma_2\sigma_3}\delta_{\sigma_3\sigma_1}\right)$, corresponding to $G' = G$,
    with all edges present, so $e(G') = 3$, $k(G') = 1$. For this term to contribute, all of the $\sigma$'s must be equal,
    amounting to $q$ possibilities, and yielding the term $qv^3$.
\end{itemize}
So combining all of these terms, one gets $Z(C_3,q,v) = q^3+3q^2v+3qv^2+qv^3$, which is precisely
$\sum_{G'\subseteq C_3}q^{k(G')}v^{e(G')}$. Now
\begin{equation}
  q^3+3q^2v+3qv^2+qv^3 = q^3+3q^2v+3qv^2+v^3+qv^3-v^3 = (q+v)^3+(q-1)v^3,
\end{equation}
so $Z(C_3,q,v) = (q+v)^3+(q-1)v^3$. By the same methods, one obtains the generalization
\begin{equation}
  Z(C_n,q,v) = (q+v)^n+(q-1)v^n .
\end{equation}
We derived this inductively in the last lecture using transfer-matrix methods, where we showed that the $q\times q$ transfer
matrix on a 1D lattice with PBC has one (leading) eigenvalue $\lambda_1 = q+v$ and a subleading eigenvalue $\lambda_2 = v$ with
multiplicity $q-1$, so $Z(C_n,q,v) = \lambda_1^n+(q-1)\lambda_2^n = (q+v)^n+(q-1)v^n$. In that derivation, $q$ had to be a
positive integer. Our present derivation does not require $q$ to be a positive integer and hence is more general.

\subsection{Chromatic and Tutte polynomials}

Next, let us define two important graph-theoretic polynomials.

\paragraph{Def.} The \emph{chromatic polynomial} $P(G,q)$ is the number of ways of assigning $q$ colors to the vertices of $G$
such that no two adjacent vertices have the same color. This is called a proper $q$-coloring of (the vertices of) $G$. The
minimum number of colors needed for a proper $q$-coloring of $G$ is the chromatic number, $\chi(G)$.

Note that if a graph contains any loop, then $P(G,q)$ vanishes identically, since a vertex with a loop is adjacent to itself,
so it cannot satisfy the constraint that it has a different color than adjacent vertices.

In the antiferromagnetic Potts model, consider the limit $T\to 0$, so $K\to-\infty$ (since $K = \beta J$ and $J<0$) and
$v = e^K-1\to-1$. In this limit, the only spin configurations that contribute to $Z(G,q,v=-1)$ are those for which adjacent
spins have different values. Hence, the $T = 0$ limit of the Potts antiferromagnet partition function is the chromatic
polynomial:
\begin{equation}
  Z(G,q,-1) = P(G,q).
\end{equation}
Since $Z(G,q,v)$ always contains an overall factor of $q$, this is also true for $P(G,q)$. Given the meaning of $P(G,q)$ as the
number of proper $q$-colorings of $G$, the fact that $P(G,0) = 0$ is clear.

There is a close connection between the Potts model partition function and a function of major importance in mathematical
graph theory, the Tutte (also called Tutte--Whitney) polynomial, $T(G,x,y)$.

\paragraph{Def.} The \emph{Tutte polynomial} of a graph $G$ is
\begin{equation}
  T(G,x,y) = \sum_{G'\subseteq G}(x-1)^{k(G')-k(G)}(y-1)^{c(G')}
           = \sum_{G'\subseteq G}(x-1)^{k(G')-k(G)}(y-1)^{e(G')+k(G')-n(G')},
\end{equation}
where $G'$ is a spanning subgraph of $G$ and $k(G')$ and $c(G')$ denote the number of components and (linearly independent)
cycles in $G'$. Since $k(G')-k(G)\ge 0$ and $c(G')\ge 0$, it follows that $T(G,x,y)$ is a polynomial in $x$ and $y$.

To prove this connection (essentially an equivalence with appropriate changes of variables), let
\begin{equation}
  x = 1+\frac qv,\qquad y = v+1,
\end{equation}
so $q = (x-1)(y-1)$ and $v = y-1$. Then, using the fact that $n(G') = n(G)$, so
$c(G') = e(G')+k(G')-n(G') = e(G')+k(G')-n(G)$, and $x-1 = q/v$, $y-1 = v$, we have
\begin{align}
  T(G,x,y) &= (x-1)^{-k(G)}(y-1)^{-n(G)}\sum_{G'\subseteq G}(x-1)^{k(G')}(y-1)^{e(G')+k(G')}\nonumber\\
  &= (x-1)^{-k(G)}(y-1)^{-n(G)}\sum_{G'\subseteq G}\left(\frac qv\right)^{k(G')}v^{e(G')+k(G')}\nonumber\\
  &= (x-1)^{-k(G)}(y-1)^{-n(G)}\sum_{G'\subseteq G}q^{k(G')}v^{e(G')}\nonumber\\
  &= (x-1)^{-k(G)}(y-1)^{-n(G)}Z(G,q,v),
\end{align}
i.e.
\begin{equation}
  Z(G,q,v) = (x-1)^{k(G)}(y-1)^{n(G)}T(G,x,y) = \left(\frac qv\right)^{k(G)}v^{n(G)}T(G,x,y).
\end{equation}

\subsubsection*{Special valuations of the Tutte polynomial}

Special valuations of $T(G,x,y)$ count various types of subgraphs of $G$. Setting $x = 1$ in $T(G,x,y)$ restricts the $G'$s
that contribute to those that are connected, i.e.\ $k(G') = k(G) = 1$. Setting $y = 1$ restricts the $G'$s that contribute to
those having no cycles, so that $c(G') = 0$.

\paragraph{Def.} A connected spanning subgraph that is a tree (is cycle-free, $c(G') = 0$) is called a \emph{spanning tree}.
Setting $x = y = 1$ picks out $G'$s that are connected spanning subgraphs with $c(G') = 0$, i.e.\ spanning trees. So
\begin{equation}
  T(G,1,1) = N_{ST}(G).
\end{equation}

\paragraph{Def.} A spanning subgraph $G'$ that has no cycles but may be disconnected, i.e.\ consist of more than one component,
is called a \emph{spanning forest} of $G$. One can relax the restriction on connectedness by setting $x = 2$, so that
$(x-1)^{k(G')-k(G)} = 1$, independent of $k(G')$. Then
\begin{equation}
  T(G,2,1) = N_{SF}(G),
\end{equation}
where $N_{SF}(G)$ is the number of spanning forests of $G$.

One can keep the restriction on connectedness but include $G'$s with cycles by setting $y = 2$; then $(y-1)^{c(G')} = 1$
independent of $c(G')$. Hence
\begin{equation}
  T(G,1,2) = N_{CSSG}(G),
\end{equation}
where $N_{CSSG}$ is the number of connected spanning subgraphs of $G$. If one sets $x = y = 2$, then the summand is just 1, so
this counts all of the spanning subgraphs of $G$: $T(G,2,2) = 2^{e(G)}$.

Special cases of the Tutte polynomial yield many graph-theoretic functions of interest. One of these is the chromatic
polynomial; for $v = -1$, i.e.\ $y = 0$, since $Z(G,q,v=-1) = P(G,q)$, one has
\begin{equation}
  P(G,q) = (-q)^{k(G)}(-1)^{n(G)}\,T(G,x=1-q,y=0).
\end{equation}

\subsection{Coloring problems and the four-color theorem}

The chromatic polynomial is of interest not just in graph theory but also in applied mathematics and engineering. One
application is the following frequency allocation problem. Consider $n$ radio broadcast transmitter stations; let each of these
be represented by a vertex of a graph $G$ and define two transmitters as being adjacent (joined by an edge) if they are closer
than a certain distance to each other. Then assign frequencies from a set of $q$ values to these transmitters, subject to the
condition that adjacent stations should use different frequencies to avoid interference. The number of ways of doing this is
$P(G,q)$.

The early interest by Birkhoff and Whitney in chromatic polynomials was motivated by their connection with the map coloring
problem, i.e.\ face-colorings of graphs; further work was done by Tutte, Read, and others.

\paragraph{Def.} A \emph{proper $q$-coloring of the faces} of a graph is an assignment of colors, from a set of $q$ colors, to the
faces of a graph such that no two adjacent faces (faces that share an edge) have the same color.

As we noted, for a planar graph $G$, there is a 1--1 correspondence between the vertices of $G$ and the faces of the planar dual
graph $G^*$. Therefore, $P(G,q)$ equivalently counts the proper $q$-colorings of the vertices of $G$ and the proper
$q$-colorings of the faces of $G^*$.

Recall that an edge of a graph such that its deletion would increase the number of components by 1 is called a bridge. A
bridge in a planar graph leads to a face being adjacent to itself across the bridge. Since adjacent faces must have different
colors in a proper $q$-coloring of the faces of a planar graph $G^*$, this is not possible if $G^*$ contains any bridges. This is
the planar dual of the property that a proper $q$-coloring of a graph $G$ is not possible if the graph contains any loops.

The \emph{Four-Color Theorem} is the statement that if $G$ is a planar graph with no loops, so $G^*$ is a planar graph with no
bridges, then $P(G,4)>0$, i.e.\ there exists a proper $q$-coloring of the faces of $G^*$ with $q = 4$ colors; equivalently,
there exists a proper $q$-coloring of the vertices of $G$ with $q = 4$ colors. Mathematicians sought to prove this for over a
century. It was proved in 1976 by Appel and Haken.

\subsection{Simple examples of \texorpdfstring{$Z(G,q,v)$, $T(G,x,y)$ and $P(G,q)$}{Z, T and P}}

We have shown that for a circuit graph, $C_n$, $Z(C_n,q,v) = (q+v)^n+(q-1)v^n$. Using $x-1 = q/v$ and $q = (x-1)(y-1)$, we
can write this as
\begin{equation}
  Z(C_n,q,v) = v^n\left[\left(1+\frac qv\right)^n+(q-1)\right] = v^n\left[x^n+(q-1)\right]
  = \left(\frac qv\right)v^n\left[\frac{x^n+(q-1)}{x-1}\right],
\end{equation}
and hence we have
\begin{equation}
  T(C_n,x,y) = \left[\frac{x^n+(xy-y-x)}{x-1}\right].
\end{equation}
Using $(x^n-x)/(x-1) = \sum_{j=1}^{n-1}x^j$ (for $n\ge 2$), one can write this as
\begin{equation}
  T(C_n,x,y) = \left(\sum_{j=1}^{n-1}x^j\right)+y .
\end{equation}
As is the case here, in expressions for $T(G,x,y)$ with a $1/(x-1)$ prefactor, it is always cancelled, so that $T(G,x,y)$ is a
polynomial in $x$ and $y$.

Setting $v = -1$ in the results for $Z(G,q,v)$ yields $P(G,q)$. For example, from $Z(C_n,q,v) = (q+v)^n+(q-1)v^n$, it follows
that
\begin{equation}
  P(C_n,q) = (q-1)^n+(q-1)(-1)^n .
\end{equation}

Many properties of $P(G,q)$ can be proved by combinatoric arguments. For example, if a graph $G$ contains any edge (i.e.\ is
not the null graph, $E_n$), then $P(G,q)$ clearly vanishes if $q = 1$, since one would have to use at least two colors to
assign different colors to the vertices that are adjacent along this edge. This, together with the property that $P(G,q)$ is a
polynomial in $q$, proves that, in addition to the factor of $q$ that it always contains, $P(G,q)$ also contains a factor of
$q-1$.

Similarly, if a graph $G$ contains any subgraph that is isomorphic to a triangle, $K_3 = C_3$, then $P(G,q)$ contains not just
the factors of $q$ and $q-1$, but also a factor of $q-2$. This is proved by noting that one can choose the color assigned to
the first vertex of the $K_3$ subgraph in any of $q$ ways, but then one can choose the color of the second vertex in $K_3$ in
only $q-1$ ways, and one can choose the color of the third vertex in the $K_3$ subgraph in only $q-2$ ways. So if $q = 2$, then
$P(G,q) = 0$. Since $P(G,q)$ is a polynomial, this proves that it contains a factor of $q-2$.

\paragraph{Def.} A graph is \emph{$k$-partite} if the vertex set $V$ can be partitioned as $V = V_1\cup V_2\cup\cdots\cup V_k$,
where $V_i\cap V_j = \emptyset$ if $i\ne j$ and no vertex in $V_i$ connects via an edge to any other vertex in $V_i$ for any $i$.

Another simple result is the following: for a connected bipartite graph $G_{\text{bip.}}$,
\begin{equation}
  P(G_{\text{bip.}},q=2) = 2 .
\end{equation}
This is proved by noting that one can choose the color to be assigned to a vertex in $V_1$ in any of 2 ways. Then one must
assign the other color to all of the vertices in $V_2$, and the first color to all the other vertices of $V_1$. As we have
discussed, the circuit graph $C_n$ is bipartite if and only if $n$ is even. By similar combinatoric arguments, one has
$P(G_{tri},3) = 3!$, where $tri$ is the triangular lattice, with free boundary conditions. To prove this, one assigns the three
colors to the vertices of one triangle, for which there are $3!$ ways. This uniquely determines the color assignments to all
other triangles on the lattice.
